%% ch08 --- The road to chaos, and a number nobody expected
%%
%% Written September 2026. Every number the reader cannot do in their head
%% comes from code/measure/ch08.py, which imports the SAME functions the
%% listings print (code/ch08/), so the page and the listings cannot disagree.
%% The logistic map's forks and landmarks use only + - * and /, so their
%% digits are the same on every machine; the sine map calls math.sin, so
%% what the page prints from it is rounded to two decimals (four for where
%% its doublings end) and checked to sit clear of a rounding boundary.
%%
%% WHAT THIS CHAPTER MAY NOT SAY. The four behaviours of the logistic map and
%% the fixed point's loss of stability at r = 3 are Chapter 5's; they are
%% recalled here, not taught. The Lyapunov exponent across r is Chapter 7's.
%% That the whole bifurcation diagram lies along a line through the
%% Mandelbrot set -- and the conversion between r and c -- is Chapter 12's
%% payoff: this chapter points at it in one sentence and computes nothing of
%% it. The flour-beetle experiment is Chapter 15's.
%%
%% Twin: chapters/pl/ch08-bifurcations.tex. Section for section, macro for
%% macro, number for number; tools/parity.py compares them.

\chapter{The road to chaos, and a number nobody expected}
\label{ch:bifurcations}
\index{bifurcation diagram}\index{Feigenbaum's constant}

\noindent
Chapter~\ref{ch:logistic} turned one knob, $r$, to four positions and found
four behaviours. At $r = \num{2.8}$ the logistic map settled on one value; at
\num{3.2} it alternated between two; at \num{3.5} it cycled through four; at
\num{3.9} it never repeated at all. Four photographs of one machine. What
happens between them? Does order turn into chaos at one setting of the knob,
the way water boils at a hundred degrees, or along a road?

The answer is one picture, among the most reproduced in science, and hidden
in it is a number nobody put there. It comes out of other equations too, it
has been measured in liquid helium, and by the end of the chapter you will
have computed it to five figures with nothing but addition, subtraction,
multiplication and division.

\begin{outcomes}
\outcome{read a bifurcation diagram and say what the map does at any value
  of the knob it shows}
\outcome{locate a period doubling by watching runs of the map, and say what
  limits how precisely that can be done}
\outcome{find a superstable parameter by bisection, using nothing but
  $+$, $-$, $\times$ and $\div$}
\outcome{compute the ratio of successive gaps between doublings, and estimate
  where the doublings run out}
\outcome{run the same search on a second rule and say what the comparison
  shows}
\outcome{find a window of order inside the chaotic range, and state what Li
  and Yorke proved about period three}
\end{outcomes}

\section{One picture for every setting of the knob}
\label{sec:ch08-picture}
\index{bifurcation diagram}\index{transient}

\noindent
To see every $r$ at once you need a recipe that turns a whole run into a few
dots, and then you draw the dots of thousands of runs side by side.

\begin{enumerate}
\item Choose a value of $r$ and start the map at $x = \num{0.5}$.
\item Run it for a thousand steps and throw them away. That stretch is the
  \emph{transient}: the orbit finding its way to wherever it is going.
\item Run it for three hundred more, and draw a dot above $r$ at the height
  of each value.
\end{enumerate}

\pyregion{code/ch08/bifurcation.py}{longrun}{The recipe: let the orbit settle, then keep what it visits}{lst:ch08-longrun}

\noindent
\code{orbit} is the book's library function that applies a rule over and
over and keeps every value; \code{xs[burn + 1:]} keeps the ones after the
transient.

\begin{think}
At $r = \num{2.8}$ the map settles on the fixed point $1 - 1/r$ that
Chapter~\ref{ch:logistic} found. Once the transient has been thrown away, how
many different dots does the recipe draw above $r = \num{2.8}$?
\end{think}
\reveal{One. All three hundred values are \val{ch08.fixed} to four decimals,
which is $1 - 1/\num{2.8}$, so the dots land on top of each other.}
A two-cycle draws two dots, a four-cycle four, and an orbit that never
repeats draws three hundred different ones, a short vertical smear. The rest
of the listing asks at six values of $r$:

\transcript{ch08-bifurcation}

\noindent
Two rows are new: at \num{3.56} the orbit cycles through eight values, and
at \num{3.567} through sixteen. Now run the recipe for thousands of values of
$r$ between \num{2.8} and $4$ and draw every dot at once.

\plotfig{ch08-diagram}{The bifurcation diagram of the logistic map. Above
each value of $r$ the dots show where $x$ spends its time once the transient
is over: one line where it settles, two where it alternates, a smear where it
never repeats.}{the bifurcation diagram}

\noindent
This is the \emph{bifurcation diagram}, the whole of
Chapter~\ref{ch:logistic} in one frame. Read it from left to right, as a
journey along the knob. One line up to $r = 3$; there it splits in two; near
\num{3.45} each branch splits again; then everything dissolves into a smear
streaked with pale stripes. A \emph{bifurcation} (from the Latin for a
two-pronged fork) is a value of the knob at which the kind of long-run
behaviour changes. Every split here is a \emph{period doubling}: each value
the orbit visited becomes two, so the cycle takes twice as long to come
round.

\begin{think}
Look at the picture between $r = \num{3.5}$, where the orbit has period
four, and the smear. Does order give way to chaos at one value of $r$, like a
switch?
\end{think}
\reveal{It looks like a switch, and it is not. Magnify the stretch before the
smear and each branch splits again, and again: periods $8$, $16$, $32$ and
on without end, all fitting in before $r = \val{ch08.end}$. The table has
already caught two of them.}

\begin{trapbox}
\textbf{\enquote{The change from order to chaos is a sudden switch.}} Not on
this road. Order breaks down in an infinite \emph{cascade} of doublings, each
arriving sooner than the one before, and chaos begins where the cascade runs
out. The smear looks sudden only because its last splits are too close
together to draw. Other systems reach chaos by other roads; this one, the
\emph{period-doubling cascade}, is the road a single hump takes.
\end{trapbox}

\section{Forks that crowd together}
\label{sec:ch08-cascade}
\index{period doubling}\index{bisection}

\noindent
Where exactly is each split? The table brackets them: period one at
\num{2.8}, period two at \num{3.2}, so the first split lies between. Try the
middle of the bracket, see which period the map settles into there, keep the
half that still contains the change, and repeat. Twenty rounds shrink the
bracket a millionfold. This is \emph{bisection}, and all it needs is a way to
ask \enquote{is it still period $p$ here?}

\pyregion{code/ch08/doublings.py}{fork}{Pinning down a split by halving the bracket}{lst:ch08-fork}

\noindent
\code{settled\_period} runs the map for twenty thousand steps and asks the
library's \code{period} function what repeats. The listing runs \code{fork}
on the four brackets the table provides.

\begin{think}
The first two splits turn out to be \val{ch08.gap.one} apart in $r$, and the
next two \val{ch08.gap.two} apart. How many times smaller is the second gap?
If the next gap shrinks by the same factor, how long will it be?
\end{think}
\reveal{About \val{ch08.ratio.one} times smaller:
$\val{ch08.gap.one} \div \val{ch08.gap.two}$. The same factor again would
make the next gap about \num{0.02}.}
Here is what the listing measured:

\transcript{ch08-doublings}

\noindent
The third gap is \val{ch08.gap.three}, and the second divided by the third is
\val{ch08.ratio.two}: a factor a little under five each time. That has a
consequence. If every gap is about \num{4.7} times the next, the gaps still
to come are the last one divided by \num{4.7}, then by \num{4.7} again, and
so on, and together they add up to the last gap divided by $\num{4.7} - 1$,
a little over a quarter of it. Infinitely many splits fit in a finite stretch
of the knob, and the cascade ends a short way beyond \val{ch08.fork.four}.

\begin{think}
The first split is known exactly. Chapter~\ref{ch:logistic} found that the
fixed point loses its grip at $r = 3$, where its slope, $2 - r$, passes
$-1$. The table says \val{ch08.fork.one}. What went wrong?
\end{think}
\reveal{Nothing in the rule: the measurement ran out of patience. At
$r = \num{2.9998}$ each step multiplies a small nudge by
$2 - r = \num{-0.9998}$, so the nudge shrinks by a fiftieth of a per cent per
step, and after twenty thousand steps the orbit is still visibly flipping
from side to side. \code{settled\_period} calls it a two-cycle.}
Every split measured by watching comes out early for the same reason: close
to a split, the cycle that is losing its grip settles very slowly.

\begin{note}
The second split can be found exactly. Once round the two-cycle, a nudge is
multiplied by the slope at each of its two points, and a page of school
algebra shows that the product is $4 + 2r - r^{2}$. That reaches $-1$ when
$r^{2} - 2r - 5 = 0$, at $r = 1 + \sqrt{6} = \val{ch08.fork.two.exact}$; the
table's \val{ch08.fork.two} is early, as expected. Later splits have no
formula this simple.
\end{note}

\section{The top of the hump}
\label{sec:ch08-superstable}
\index{superstable cycle}\index{Feigenbaum, Mitchell}

\noindent
In 1975 Mitchell Feigenbaum, a physicist at the Los Alamos laboratory in New
Mexico, was working out where these splits happen on a programmable pocket
calculator, and noticed the pattern you have just seen: each gap a fixed
fraction of the one before. To measure that fraction sharply, use a
better-placed landmark than the split, where a cycle hangs on by its
fingertips: the point in each stretch where the cycle is held most firmly.

The slope of the logistic map at $x$ is $r(1 - 2x)$, the stretch factor
Chapter~\ref{ch:feedback} taught you to measure. It is zero at
$x = \num{0.5}$, the top of the hump. A cycle through $\num{0.5}$ multiplies
any nudge by that zero once per trip, wiping it out. Such a cycle is called
\emph{superstable}, and for each period $1, 2, 4, 8, \ldots$ there is exactly
one value of $r$ on the road where it happens: this chapter calls those
values \emph{landmarks}. Each is the answer to one equation: start at
$\num{0.5}$, take as many steps as the period, and arrive back at $\num{0.5}$.

\begin{think}
For period one: for which $r$ does one step take $\num{0.5}$ straight back to
$\num{0.5}$? The step is $r \times \num{0.5} \times \num{0.5}$.
\end{think}
\reveal{$r = 2$. Half of a half is a quarter, and $2 \times \num{0.25} =
\num{0.5}$.}
For longer periods there is no mental shortcut, and bisection needs none.
\code{from\_top} reports how far from $\num{0.5}$ the orbit ends, which is
positive on one side of a landmark and negative on the other:

\pyregion{code/ch08/superstable.py}{search}{Start at the top, and halve the bracket until the orbit comes back to it}{lst:ch08-search}

\noindent
The first two landmarks are bracketed by eye from the picture. After that
each gap is smaller than the last, so the next landmark is sought between a
tenth and a half of the last gap further on, a generous bracket for gaps
shrinking by a factor near five. A \emph{family} is a function that, given
$r$, builds the rule.

\pyregion{code/ch08/superstable.py}{landmarks}{The landmarks for periods $1, 2, 4, 8, \ldots$, each bracketed from the last}{lst:ch08-landmarks}

\noindent
The listing finds the first \val{ch08.landmark.count} landmarks, for periods
$1$ to $1024$, with nothing but $+$, $-$, $\times$ and $\div$, so every digit
is the same on every computer:

\transcript{ch08-superstable}

\noindent
Look down the last column, each gap divided by the next. It wobbles for a
few rows and then stops moving: \val{ch08.delta}. That is
\emph{Feigenbaum's constant}, written $\delta$ (the Greek letter delta). The
landmarks give it far more sharply than the splits did, because at a
landmark the orbit settles at once instead of barely at all. Splits and
landmarks take turns along the road, one of each per period, so both crowd
in at the same rate towards the same end, $r = \val{ch08.end}$. The
period-two landmark, \val{ch08.super.two}, is exactly $1 + \sqrt{5}$, where
the note's product $4 + 2r - r^{2}$ is zero.

Feigenbaum published in 1978; Pierre Coullet and Charles Tresser in France
found the same number independently at about the same time. What he did
next is why the number is famous.

\begin{lab}{l08_01_superstable}{Find a landmark yourself}
Write \code{from\_top} and \code{bisect}, and put them together in
\code{superstable}. Check your answer for period two against
$1 + \sqrt{5}$. Then find the $r$ between \num{3.82} and \num{3.84} at which
$\num{0.5}$ comes back to itself after \emph{three} steps. You will meet
that number again.
\end{lab}

\section{A number that belongs to the road}
\label{sec:ch08-universal}
\index{universality}\index{sine map}

\noindent
Feigenbaum's next step was to try a different rule, the \emph{sine map}
\[ x_{n+1} = s \sin(\pi x_{n}), \]
which returns $s$ times the sine of $\pi x$. You need no trigonometry to
follow it: for $x$ between $0$ and $1$, the sine of $\pi x$ is a smooth hump
rising from $0$ to $1$ at $x = \num{0.5}$ and falling back to $0$. So the
sine map shares the logistic map's shape \dash{} one hump, its top at
$\num{0.5}$ \dash{} and nothing of its formula: no $x(1 - x)$, a different
knob $s$, a curve that is not a parabola.

\begin{think}
Run the same landmark search on the sine map. Do you expect its gaps to
shrink by the logistic map's \val{ch08.delta} again, or by some other number
that belongs to the sine?
\end{think}
\reveal{The same number. The sine map's gaps shrink by \val{ch08.sine.first}
at first and then close in, landmark by landmark, on \val{ch08.sine.delta},
where the logistic map's ratios settle too.}
The listing knows nothing about the sine map but its family and two brackets
read off its diagram:

\pyregion{code/ch08/two_maps.py}{compare}{The same search on two unrelated rules}{lst:ch08-compare}

\transcript{ch08-two-maps}

\noindent
The two roads end in different places (the last row: \val{ch08.end} and
\val{ch08.sine.end}) and start differently, yet they settle on the same ratio
to every digit printed. The function \code{ratios} divides each gap by the
next; the second lab asks you to write it.

\plotfig{ch08-ratios}{Each gap between landmarks divided by the next, for two
rules with nothing in common but a single smooth hump. They start apart and
end on the same number.}{the ratios for two rules converging}

\begin{trapbox}
\textbf{\enquote{A different equation gives a different route.}} It gives
different signposts \dash{} the splits sit at different values of the knob
\dash{} but the road has the same shape, and its splits crowd together at the
same rate. That number belongs to the route, not to the equation. This is
\emph{universality}, and it is why a constant of nature can come out of a
toy model.
\end{trapbox}

\noindent
Why? Each new split is decided by a smaller and smaller stretch of the map
near the top of its hump, seen through a stronger and stronger magnifying
glass. Magnify any smooth, rounded hump enough around its top and it looks
like a parabola; whatever made the sine map a sine has been magnified out of
sight. The constants of school geometry, like $\pi$, belong to shapes. This
one belongs to a process.

\mermaidfig{ch08-universal}{Universality in one line: any rule with one
smooth hump takes the same road, and the road carries its own
ratio.}{any hump, one route, one ratio}

\noindent
It was a shock because it predicted something about systems whose equations
nobody could write down: if one reaches chaos by period doubling, its splits
should crowd together at Feigenbaum's rate, whatever it is made of.

\begin{worldbox}
Around 1980 Albert Libchaber and Jean Maurer in Paris heated a tiny box of
liquid helium from below and recorded the temperature at one point as the
fluid churned. As they turned up the heat, the temperature's rhythm doubled
its period, and doubled again. Later experiments on convecting fluids and
electronic circuits measured the ratio of successive gaps and found numbers
close to Feigenbaum's: not to four decimals, since a laboratory can resolve
only the first few doublings, but close enough that nobody takes it for a
coincidence. In 1986 Feigenbaum and Libchaber shared the Wolf Prize in
physics.
\end{worldbox}

\begin{rigourbox}
Why every smooth hump gives the same $\delta$ is not proved here. Feigenbaum
explained it with a method from physics, \emph{renormalisation}, which
compares the map after each doubling with a rescaled copy of itself; Oscar
Lanford gave a computer-assisted proof in 1982. \emph{Smooth} matters: a hump
whose top is flatter than a parabola's gives a different constant.
\end{rigourbox}

\begin{lab}{l08_02_ratio}{Feigenbaum's ratio, and where the doublings end}
Write \code{ratios}, which divides each gap by the next, and
\code{where\_it\_ends}, which adds up every gap still to come, assuming each
is the one before divided by the last ratio. The test feeds your functions
the landmarks of both maps: one ratio, two different ends.
\end{lab}

\section{Order inside the chaos}
\label{sec:ch08-windows}
\index{periodic window}\index{Li--Yorke theorem}

\noindent
Past $r = \val{ch08.end}$ the doublings are used up, and the diagram turns
to smear.

\begin{think}
From $r = \val{ch08.end}$ to $r = 4$, is the logistic map chaotic at every
setting of the knob, the way it is at \num{3.9}?
\end{think}
\reveal{No. Look at the picture again: the smear is streaked with pale
stripes, and in each one the orbit is back on a short cycle. The widest,
near \num{3.83}, is a window of period three.}
\code{code/ch08/windows.py} walks $r$ across that stripe with the same
\code{settled\_period}, then runs the landmark search inside it:

\transcript{ch08-windows}

\noindent
At \num{3.8284} the orbit never repeats; at \num{3.8285} it has settled into
a three-cycle. The window opens between the two, at exactly
$r = 1 + \sqrt{8} = \val{ch08.window.open}$. The three-cycle then doubles to
six and twelve, and by \num{3.85} it has dissolved into smear again. The
whole road repeats in miniature. The window's first landmark is the $r$ you
found in the first lab, \val{ch08.super.three}, and its landmarks crowd
together at a ratio of \val{ch08.window.delta}: Feigenbaum's number once
more, now from periods $3, 6, 12, \ldots$ instead of $1, 2, 4, \ldots$.

\plotfig{ch08-zoom}{Left: the period-three window. Right: the window's middle
branch, magnified: a small copy of the whole bifurcation diagram, with its
own cascade, its own smear and its own windows.}{the period-three window,
magnified}

\begin{trapbox}
\textbf{\enquote{Once a system has reached chaos, it is chaotic everywhere
beyond that point.}} Between $r = \val{ch08.end}$ and $4$ the logistic map
returns to order infinitely many times, in windows of every period. Turn the
knob a little and a chaotic system can become perfectly periodic; turn it a
little further and it is chaotic again. Chapter~\ref{ch:lyapunov}'s exponent
dips below zero in exactly these stripes.
\end{trapbox}

\begin{note}
In 1975 Tien-Yien Li and James Yorke published a paper whose title was its
result: \emph{Period three implies chaos}. Paraphrased: if a continuous rule
takes an interval into itself, as the logistic map does from $0$ to $1$, and
some start comes back to itself after exactly three steps, then the rule has
cycles of every length, and starts whose orbits never settle into any cycle.
The paper gave the word \emph{chaos} its mathematical meaning; a stronger
result had been proved in 1964 by Oleksandr Sharkovsky in Kyiv, little known
in the West until later. In the window almost every start ends on the
three-cycle, but the theorem still holds: the other cycles are there,
unstable, and an orbit may wander among them for a long time before settling.
\end{note}

\noindent
The diagram has one more surprise, and it belongs to
Chapter~\ref{ch:mandelbrot}: this whole picture lies along a line through the
most famous image in mathematics, with the splits and the window exactly
where a one-line conversion puts them.

\begin{summarybox}
\begin{itemize}
\item The \textbf{bifurcation diagram} shows, for each value of the knob,
  where the orbit lives once the transient is over. A \textbf{bifurcation} is
  a value at which the long-run behaviour changes kind; on this road each is
  a \textbf{period doubling}.
\item The splits crowd together, each gap a little under five times the next,
  so infinitely many fit before $r = \val{ch08.end}$. Measured by watching,
  a split comes out early, because near it the orbit settles very slowly.
\item At a \textbf{superstable} landmark the cycle passes through the top of
  the hump, where the slope is zero. Each is one equation, solved by
  \textbf{bisection} with $+$, $-$, $\times$ and $\div$, and the ratio of
  their gaps settles on \textbf{Feigenbaum's constant}, \val{ch08.delta}.
\item The sine map gives the same ratio: the number belongs to the route
  every smooth one-hump rule takes, not to a formula. That is
  \textbf{universality}, and it has been measured in fluids and circuits.
\item Beyond the cascade, chaos is interrupted by \textbf{windows} of order.
  The widest, of period three, opens at $1 + \sqrt{8}$ and holds a miniature
  of the whole road; Li and Yorke proved that period three implies chaos.
\end{itemize}
\end{summarybox}

\canyou

\begin{checkpoint}
\item What does the column of dots above one value of $r$ in a bifurcation
  diagram show, and what does it look like for a fixed point, a two-cycle and
  chaos?
  \answerto{Where the orbit lives in the long run at that $r$, after the
  transient: one dot for a fixed point, two for a two-cycle, a smear for
  chaos.}
\item Two successive gaps between splits are \num{0.08} and \num{0.017}.
  Roughly what is their ratio, and how long should the next gap be?
  \answerto{About \num{4.7}, from $\num{0.08} \div \num{0.017}$; so the next
  gap should be about $\num{0.017} \div \num{4.7}$, roughly \num{0.0036}.}
\item Why does a split located by watching the orbit come out slightly early?
  \answerto{Near a split the cycle losing its grip multiplies a nudge by
  nearly $-1$ each time round, so it settles extremely slowly, and a finite
  run still shows it flipping between values.}
\item What is a superstable landmark, and why can it be found to many digits
  with $+$, $-$, $\times$ and $\div$ alone?
  \answerto{The $r$ at which the cycle passes through the top of the hump,
  $x = \num{0.5}$, where the slope is zero. It solves one equation (start at
  $\num{0.5}$, take the period's number of steps, be back at $\num{0.5}$),
  which bisection solves by halving a bracket.}
\item The sine map's doublings run out at about \val{ch08.sine.end}, the
  logistic map's at about \val{ch08.end}. Does that contradict universality?
  \answerto{No. Universality says the ratio of successive gaps tends to the
  same number, about \val{ch08.delta}, for every rule with one smooth hump;
  where each road ends depends on the rule.}
\item Is the logistic map chaotic at $r = \num{3.83}$? What did Li and Yorke
  prove about period three?
  \answerto{No: \num{3.83} lies in the period-three window, which opens at
  $1 + \sqrt{8}$, about \val{ch08.window.open}. Li and Yorke proved that a
  continuous rule taking an interval into itself, with a point of period
  three, has cycles of every length and starts that never settle into any
  cycle.}
\end{checkpoint}
